% GENERATED FILE. Edit canonical lessons, then run npm run generate:books.
% canonical-source-hash: fnv1a64:82386f4af74a147e
% canonical-lessons: KA-C276-bai, KA-C276-regisu, KA-C276-jagalavadu, KA-C276-samadhana-madu, KA-C276-hemmepadu

\chapter{To scold, To tease, To quarrel, To console, To be proud}
\label{ch:ka-to-scold-to-tease-to-quarrel-to-console-to-be-proud}
\hlchaptermodality{276}

\begin{chapteropening}
\textbf{By the end of this chapter:} \emph{I can say to scold, to tease, to quarrel, to console and to be proud.}

Complete the last lesson of chapter 276: I can say to scold, to tease, to quarrel, to console and to be proud.
\end{chapteropening}

\section[\texorpdfstring{\kn{ಬೈ} (bai)}{ಬೈ (bai)}]{\kn{ಬೈ} (bai) — to scold}

\label{lesson:KA-C276-bai}



\begin{quote}
Before the new one: say the Kannada for to apply, then the Kannada for to compare.
\end{quote}

\subsection*{You'll want to know: \kn{ಬೈ}}
\textbf{\kn{ಬೈ}} — \emph{bai} — \textquotedblleft{}to scold\textquotedblright{}.

\begin{grammarlens}[title={What you've built}]
One more word for this chapter.
\end{grammarlens}

\subsection*{Your turn}
\begin{itemize}
\raggedright
  \item \emph{Say it:} \emph{bai}
  \item \emph{Say it:} \emph{bai}, once more
  \item \emph{Say it:} say \emph{bai}
  \item \emph{Recall:} say the Kannada for to apply, then the Kannada for to compare, then say \emph{bai} again
\end{itemize}

\begin{tcolorbox}[breakable,colback=teal!4,colframe=teal!35!black,title={Before you move on}]
What does \kn{ಬೈ} mean? (\textquotedblleft{}To scold\textquotedblright{}.) Say it once more.
\end{tcolorbox}
\section[\texorpdfstring{\kn{ರೇಗಿಸು} (rēgisu)}{ರೇಗಿಸು (rēgisu)}]{\kn{ರೇಗಿಸು} (rēgisu) — to tease}

\label{lesson:KA-C276-regisu}



\begin{quote}
Before the new one: say the Kannada for to compare, then the Kannada for to scold.
\end{quote}

\subsection*{You'll want to know: \kn{ರೇಗಿಸು}}
\textbf{\kn{ರೇಗಿಸು}} — \emph{rēgisu} — \textquotedblleft{}to tease\textquotedblright{}.

\begin{grammarlens}[title={What you've built}]
One more word for this chapter.
\end{grammarlens}

\subsection*{Your turn}
\begin{itemize}
\raggedright
  \item \emph{Say it:} \emph{rēgisu}
  \item \emph{Say it:} \emph{rēgisu}, once more
  \item \emph{Say it:} say \emph{rēgisu}
  \item \emph{Recall:} say the Kannada for to compare, then the Kannada for to scold, then say \emph{rēgisu} again
\end{itemize}

\begin{tcolorbox}[breakable,colback=teal!4,colframe=teal!35!black,title={Before you move on}]
What does \kn{ರೇಗಿಸು} mean? (\textquotedblleft{}To tease\textquotedblright{}.) Say it once more.
\end{tcolorbox}
\section[\texorpdfstring{\kn{ಜಗಳವಾಡು} (jagaḷavāḍu)}{ಜಗಳವಾಡು (jagaḷavāḍu)}]{\kn{ಜಗಳವಾಡು} (jagaḷavāḍu) — to quarrel}

\label{lesson:KA-C276-jagalavadu}



\begin{quote}
Before the new one: say the Kannada for to scold, then the Kannada for to tease.
\end{quote}

\subsection*{You'll want to know: \kn{ಜಗಳವಾಡು}}
\textbf{\kn{ಜಗಳವಾಡು}} — \emph{jagaḷavāḍu} — \textquotedblleft{}to quarrel\textquotedblright{}.

\begin{grammarlens}[title={What you've built}]
One more word for this chapter.
\end{grammarlens}

\subsection*{Your turn}
\begin{itemize}
\raggedright
  \item \emph{Say it:} \emph{jagaḷavāḍu}
  \item \emph{Say it:} \emph{jagaḷavāḍu}, once more
  \item \emph{Say it:} say \emph{jagaḷavāḍu}
  \item \emph{Recall:} say the Kannada for to scold, then the Kannada for to tease, then say \emph{jagaḷavāḍu} again
\end{itemize}

\begin{tcolorbox}[breakable,colback=teal!4,colframe=teal!35!black,title={Before you move on}]
What does \kn{ಜಗಳವಾಡು} mean? (\textquotedblleft{}To quarrel\textquotedblright{}.) Say it once more.
\end{tcolorbox}
\section[\texorpdfstring{\kn{ಸಮಾಧಾನ} \kn{ಮಾಡು} (samādhāna māḍu)}{ಸಮಾಧಾನ ಮಾಡು (samādhāna māḍu)}]{\kn{ಸಮಾಧಾನ} \kn{ಮಾಡು} (samādhāna māḍu) — to console, to calm (someone) down}

\label{lesson:KA-C276-samadhana-madu}



\begin{quote}
Before the new one: say the Kannada for to tease, then the Kannada for to quarrel.
\end{quote}

\subsection*{You'll want to know: \kn{ಸಮಾಧಾನ} \kn{ಮಾಡು}}
\textbf{\kn{ಸಮಾಧಾನ} \kn{ಮಾಡು}} — \emph{samādhāna māḍu} — \textquotedblleft{}to console, to calm (someone) down\textquotedblright{}.

\begin{grammarlens}[title={What you've built}]
One more word for this chapter.
\end{grammarlens}

\subsection*{Your turn}
\begin{itemize}
\raggedright
  \item \emph{Say it:} \emph{samādhāna māḍu}
  \item \emph{Say it:} \emph{samādhāna māḍu}, once more
  \item \emph{Say it:} say \emph{samādhāna māḍu}
  \item \emph{Recall:} say the Kannada for to tease, then the Kannada for to quarrel, then say \emph{samādhāna māḍu} again
\end{itemize}

\begin{tcolorbox}[breakable,colback=teal!4,colframe=teal!35!black,title={Before you move on}]
What does \kn{ಸಮಾಧಾನ} \kn{ಮಾಡು} mean? (\textquotedblleft{}To console, to calm (someone) down\textquotedblright{}.) Say it once more.
\end{tcolorbox}
\section[\texorpdfstring{\kn{ಹೆಮ್ಮೆಪಡು} (hemmepaḍu)}{ಹೆಮ್ಮೆಪಡು (hemmepaḍu)}]{\kn{ಹೆಮ್ಮೆಪಡು} (hemmepaḍu) — to be proud}

\label{lesson:KA-C276-hemmepadu}



\begin{quote}
Before the new one: say the Kannada for to quarrel, then the Kannada for to console.
\end{quote}

\subsection*{You'll want to know: \kn{ಹೆಮ್ಮೆಪಡು}}
\textbf{\kn{ಹೆಮ್ಮೆಪಡು}} — \emph{hemmepaḍu} — \textquotedblleft{}to be proud\textquotedblright{}.

\begin{grammarlens}[title={What you've built}]
One more word for this chapter.
\end{grammarlens}

\subsection*{Your turn}
\begin{itemize}
\raggedright
  \item \emph{Say it:} \emph{hemmepaḍu}
  \item \emph{Say it:} \emph{hemmepaḍu}, once more
  \item \emph{Say it:} say \emph{hemmepaḍu}
  \item \emph{Recall:} say the Kannada for to quarrel, then the Kannada for to console, then say \emph{hemmepaḍu} again
\end{itemize}

\begin{tcolorbox}[breakable,colback=teal!4,colframe=teal!35!black,title={Before you move on}]
What does \kn{ಹೆಮ್ಮೆಪಡು} mean? (\textquotedblleft{}To be proud\textquotedblright{}.) Say it once more.
\end{tcolorbox}
